\section{Examples: A Catalogue of Verified Presentations}\label{sec:examples}

This section is the library's harvest.  Table~\ref{tab:presentations} lists
every concrete \aid{{\_}IsPresentationOf{\_}} theorem in the development;
Table~\ref{tab:nfs} lists the normal-form and uniqueness witnesses behind
them.  All of them are restated with full types in the index module
\texttt{MainTheorems.agda}, which also re-exports the conditional
construction theorems of \S\ref{sec:constructions} with their premise
telescopes.  We then discuss each family, ending with the three
amalgamation case studies, whose final theorems are monoid
\emph{isomorphisms} between gate-set presentations and amalgamated
products.

\begin{table}
\caption{Concrete presentation theorems (all verified; names as in
\texttt{MainTheorems.agda}).  $p$ ranges over odd primes, $n, m, N$ over
naturals.}
\label{tab:presentations}
\centering
\small
\begin{tabular}{@{}lll@{}}
\toprule
Relation set & Presents & Index name \\
\midrule
\aid{EmptyRel} over $\bot$ & trivial group & \aid{trivial-presentation} \\
\aid{TrivialRel} over any $A$ & trivial group & \aid{trivial-presentation′} \\
$T^{\,n+1} = \varepsilon$ & $\mathbb{Z}/(n{+}1)\mathbb{Z}$ & \aid{cyclic-presentation} \\
\aid{order}, \aid{yang-baxter} + structural & permutations of $\mathsf{Fin}\;n$ & \aid{symmetric-presentation} \\
$C_p \aid{⋄} C_p \aid{⋄} \aid{CommRel}$ & $\mathbb{Z}_p \times \mathbb{Z}_p$ & (\aid{H-pres}, Pauli development) \\
$(C_p \aid{⋄} C_p \aid{⋄} \aid{CommRel})\,\aid{⊕\char94}\,n$ & Pauli quotient $(\mathbb{Z}_p{\times}\mathbb{Z}_p)^n$ & \aid{pauli-presentation} \\
$(C_N\,\aid{⊕\char94}\,n) \aid{⋄} S_n\text{-circuits} \aid{⋄} \aid{ConjRelʷ}$ & $\mathbb{Z}_N \wr S_n$ (signed permutations) & \aid{wreath-presentation} \\
\bottomrule
\end{tabular}
\end{table}

\begin{table}
\caption{Normal-form and uniqueness witnesses (selection).  NFI =
\aid{NormalFormInjective}, NF = \aid{NormalForm} (with section), UNF =
\aid{UniqueNormalForm}.}
\label{tab:nfs}
\centering
\small
\begin{tabular}{@{}llll@{}}
\toprule
Presentation & Carrier & Witnesses & Note \\
\midrule
cyclic, order $n{+}1$ & $\mathsf{Fin}\,(n{+}1)$ & NFI, NF, UNF & order $0$ gives $\mathbb{Z}$, carrier $\mathbb{N}$ \\
$S_n$ circuits & $\top \times C\,1 \times \dots \times C\,(n{-}1)$ & NFI, NF, UNF$\times 2$ & factorial system, $|{\cdot}| = n!$ \\
$\Gamma \oplus \Delta$ & $\mathit{NF}_1 \times \mathit{NF}_2$ & NFI, NF & lifted from factors \\
$\Gamma \rtimes \Delta \star \mathit{conj}$ & $\mathit{NF}_1 \times \mathit{NF}_2$ & NFI, NF & twist on the left factor \\
amalgamated $P_1 \ast P_2$ & $(D \uplus \top) \times \mathsf{List}(C {\times} D) \times (C \uplus \top)$ & NFI, NF & alternating cosets \\
Clifford+$T$ base $\langle \omega \rangle {\times} \langle S \rangle$ & $\mathsf{Fin}\,8 \times \mathsf{Fin}\,4$ & NFI, NF & radices $8, 4$ \\
\bottomrule
\end{tabular}
\end{table}

\subsection{Cyclic groups, and the trivial group twice}

The cyclic development presents $\langle\, T \mid T^{\,n+1} =
\varepsilon\,\rangle$ with normal-form carrier $\mathsf{Fin}\,(n{+}1)$ ---
and, pleasingly, the \emph{same} relation family at index $0$ presents the
infinite cyclic group $\mathbb{Z}$ with carrier $\mathbb{N}$ (a word count
with no wraparound; the presented monoid over one generator with no
relation is $\mathbb{N}$, which is already a group presentation of nothing
--- grouplikeness fails --- so the library states the group results for
positive orders and keeps the monoid-level normal form uniformly).  The
trivial group is presented twice --- empty alphabet with no relations, and
arbitrary alphabet with the everything-is-$\varepsilon$ relation --- and
the two presentations are proved isomorphic; besides being the base case
of $n$-fold products, this pair is a two-page worked example of the whole
pipeline and a good entry point into the code.

\subsection{A product demonstration: $S_3 \times C_5$}

A small worked example composes two unrelated developments: the direct
product of the $3$-wire permutation presentation and the cyclic
presentation of order $5$, with the lifted normal form of
\S\ref{sec:constructions} on the carrier $\mathit{NF}_{S_3} \times
\mathsf{Fin}\,5$.  The example exists to be read: it shows the transport
machinery end to end on a group of order $30$, with the monoid-level
semantics assembled from the factor semantics.  (A presentation-level
upgrade of this demonstration is work in progress and not claimed here.)

\subsection{Wreath products: signed permutations and beyond}

The headline compositional result combines almost every construction in
the library.  Fix a base order $N = m{+}1$ and $n$ wires.  The $n$-fold
direct power $(C_N)^{\oplus n}$ presents $(\mathbb{Z}_N)^n$; the circuit
presentation of $S_n$ acts on it by permuting coordinates --- an action
specified \emph{syntactically} as a word-valued conjugation
$\mathit{conj} : \aid{Gen}\;n \to (\top \uplus^n) \to
\aid{Word}\,(\top \uplus^n)$ sending a generator and a coordinate to the
permuted coordinate.  Two finite compatibility lemmas (the action
respects the Coxeter relations pointwise, and respects the product
relations) feed the semidirect-product presentation theorem, giving
\begin{agdacode}
wreath-presentation : ∀ n m →
  ((((suc m) Cn,_===_) ⊕^ n) ⋄ (n VRel,_===_) ⋄ ConjRelʷ (SnD.conj {n}))
    IsPresentationOf (SnD.Wreath.wreath-group n m)
\end{agdacode}
where $\aid{wreath-group}\;n\;m$ is the semantic
$(\mathbb{Z}_N)^n \rtimes S_n = \mathbb{Z}_N \wr S_n$ built from the
\texttt{ForStdlib} semidirect product and the bundled permutation group.
At $N = 2$ this is the group of \emph{signed permutations} --- the
hyperoctahedral group $B_n$, the symmetry group of the hypercube; at
general $N$ it is the generalized symmetric group $G(N,1,n)$ of monomial
matrices with $N$-th-root-of-unity entries.  These groups are ubiquitous
in the qudit setting (Pauli-frame changes, permutation-and-phase
subgroups of Clifford groups), and the theorem delivers their
presentations \emph{for all $N$ and $n$ at once} --- a family of
completeness results parameterized over two naturals, which is where a
mechanized, compositional treatment clearly outruns the per-instance
pen-and-paper approach.

\subsection{Pauli groups in all odd prime dimensions}\label{sec:pauli}

The Pauli development is the purest showcase of compositionality --- there
is \emph{no} coset enumeration in it at all:

\begin{agdacode}
pauli-presentation :
  ∀ (p-2 : ℕ) (p-prime : Prime (2+ p-2)) (n : ℕ) →
  (Pauli.Γ-H p-2 p-prime ⊕^ n) IsPresentationOf (Pauli.Pauli-group p-2 p-prime n)
\end{agdacode}

For an odd prime $p$, the single-qupit relation set is the direct-product
join of two order-$p$ cyclic presentations ($X$- and $Z$-type generators,
which commute in the quotient of the Pauli group by its phase subgroup),
and the $n$-qupit relation set is its $n$-fold power.  The semantic group
$(\mathbb{Z}_p \times \mathbb{Z}_p)^n$ is assembled by the standard
library's binary direct product and our $n$-fold iteration, and the
presentation is literally three construction theorems composed.  The
full Pauli group with phases, and its role beneath the Clifford
hierarchy, belongs to the symplectic work in progress
(\S\ref{sec:wip}).

\subsection{Single-qubit Clifford+$T$, as an amalgamated free product}\label{sec:cliffordt1}

The first amalgamation case study presents the single-qubit Clifford+$T$
\emph{monoid} --- the flagship example of the completeness literature
--- by the relation set (over generators $T, H, S, \omega$, with
$X := H S S H$ as an abbreviation):
\[
  \omega^8 = \varepsilon,\quad
  T^2 = S,\quad
  S^4 = \varepsilon,\quad
  H^2 = \varepsilon,\quad
  (SH)^3 = \omega,\quad
  (TX)^2 = \omega,\quad
  \omega\ \text{central}.
\]
The development builds a tower whose base is the direct product
$\langle \omega \rangle \times \langle S \rangle$ of cyclic groups of
orders $8$ and $4$ (normal-form radices $8 \cdot 4$), extends it by one
verified coset table to the subgroup
$\mathit{XS\omega} = \langle X, S, \omega \rangle$ (two cosets:
$\varepsilon$ and $X$), and then extends \emph{twice from the same base}:
by a two-coset table to
$\mathit{TXS\omega} = \langle T, X, S, \omega \rangle$, and by a
three-coset table (cosets $\varepsilon, H, HS$) to the Clifford group
$\langle H, X, S, \omega \rangle$.  The two \aid{PackedCosetTable}s over
the shared $\mathit{XS\omega}$ form an \aid{AmalDataNF}, and the
amalgamation module produces the normal form for the amalgamated product
--- alternating $T$-side and $H$-side coset representatives around a
common-subgroup tail.  It is instructive to see what this carrier
\emph{is}: with a one-element $C$ ($T$-coset) and two-element $D$
($H, HS$), an element of $\aid{amalNFC}\;C\;D$ is an optional leading
$H$-syllable, a list of alternations, and an optional trailing $T$ ---
precisely the Matsumoto--Amano normal form
$(\varepsilon \mid T)\,(HT \mid SHT)^{*}\,\mathcal{C}$ read
group-theoretically~\cite{matsumoto2008representation,giles2013remarks}:
the alternating normal form of an amalgamated free product, as observed
for this group in the Bian--Selinger
line~\cite{bian2023cliffordt2,bian2023thesis}.  The final theorem is the
syntactic isomorphism assembled by \aid{StarIsomorphism} from
generator-level translations both ways:

\begin{agdacode}
clifford+T-qubit-isomorphism :
  MonoidMorphisms.IsMonoidIsomorphism
    (Monoid.rawMonoid (PP.•-ε-monoid CliffordT1.CliffordT1._===_))
    (Monoid.rawMonoid (PP.•-ε-monoid CliffordT1.CliffordT1.amalPres))
    (CliffordT1.CliffordT1.f ʷ)
\end{agdacode}

--- the flat gate-set presentation and the amalgamated-product
presentation define isomorphic monoids, and both carry the verified
normal form.  Every coset-table obligation in the tower (well-definedness
on all axioms at all cosets --- several dozen cases) is discharged by
\aid{by-equal-nf Eq.refl}: the typechecker computes both coset
descriptors and observes they are equal.

\subsection{Single-qutrit Clifford+$T$}

The qutrit analogue replaces $\omega$ by a ninth root $\zeta$ and runs a
longer tower: from the cyclic base $\langle \zeta \rangle$,
$\zeta^9 = \varepsilon$, through
$\langle S, \zeta \rangle$ ($S^3 = \zeta^6$; three cosets), then a
\emph{nine}-coset table to $\langle S, X, \zeta \rangle$
($X^3 = \varepsilon$, $(SX)^3 = \varepsilon$, and a twisted commutation
$(XS)(SX) = \zeta^6 (SX)(XS)$), then a two-coset table adjoining the
squared Hadamard $H\!H$, and finally an amalgamation of a $T$-side and an
$H$-side extension over the common subgroup.  The presented gate-set
relations (generators $T, H, S, \zeta$; abbreviations
$X := \zeta^3 (HSH)(HSSH)$ and $Z := \zeta^3 S^2 H^2 S H^2$):
\[
  \zeta^9 = \varepsilon,\;\;
  S^3 = \zeta^6,\;\;
  H^4 = \varepsilon,\;\;
  T^3 = Z,\;\;
  (SH)^3 = \varepsilon,\;\;
  (THH)^2 = \varepsilon,
\]
\[
  HHSHHS = SHHSHH,\;\;
  TS = ST,\;\;
  TX = \zeta^3 S^2 X\, T,\;\;
  \zeta\ \text{central},
\]
a presentation derived and verified in this development (its normal form
refines the single-qutrit and single-qudit Clifford+$T$ normal forms
studied on the matrix side~\cite{prakash2018qutrit,jain2020qudit}).  The final theorem has the same shape as
the qubit one --- \aid{clifford+T-qutrit-isomorphism} in the index ---
and the same proof economy: the nine-coset level alone imposes
seventy-two axiom-preservation obligations (nine cosets against eight
axiom cases), every one closed by \aid{refl} after evaluation.  We know of no prior machine-checked completeness
result for a qutrit gate-set presentation.

\subsection{$U_3(\mathbb{Z}[\tfrac12,i])$, a two-level amalgamation}

The third case study formalizes the case $n = 3$ of the presentation, due
to Bian and Selinger, of the groups
$U_n(\mathbb{Z}[\tfrac12,i])$~\cite{bian2021un} --- unitary matrices over
the ring $\mathbb{Z}[\tfrac12,i]$, the natural home of the
Clifford+$T$-with-ancilla gate set
$\{X, CX, CCX, \omega^\dagger H, S\}$~\cite{amy2020numbertheoretic}.
Generators are $i_0$ (a local phase), $K_{01}$ (a $2{\times}2$
Hadamard-type block), and the two transpositions $X_{01}, X_{12}$;
derived generators $i_1, i_2, K_{02}, K_{12}$ are conjugates.  The eleven
relations ($\mathrm{S1}$--$\mathrm{S14}$ in the development, following the
paper's numbering) include the orders $i_0^4 = X_{01}^2 = X_{12}^2 =
\varepsilon$, braid and commutation laws for the transpositions and
phases, and the genuinely quantum interactions of $K_{01}$ with the
phases, e.g.
$K_{01} i_1^2 = X_{01} K_{01}$ and
$K_{01}^2\, i_0\, i_1 = \varepsilon$.  The verified structure is a
\emph{two-level} amalgamation: a six-coset Reidemeister--Schreier level
presents the $\langle K_{01}, i_0, i_1 \rangle$ subgroup over the product
of two order-$4$ cyclic groups, definitional extensions
(\aid{SugarRel}) introduce the derived generators, and an amalgamated
product glues the permutation part (built from the $S_3$ development of
\S\ref{sec:permutations}) to the $K$-part over their common phase
subgroup.  The final theorem,
\aid{U₃-presentation-isomorphism}, is again a monoid isomorphism between
the flat eleven-relation presentation and the amalgamated-product
presentation, with the alternating normal form on the amalgam side.  At
$2\,869$ and $1\,395$ source lines, the qutrit and $U_3$ studies are the
largest verified objects in the library, and both consist mostly of
first-order tables and \aid{refl}-discharged obligations --- the shape of
proof this framework was designed to make routine.

\subsection{The qubit Clifford groups as an extension, with an explicit cocycle}\label{sec:cliffordqubit}

The catalogue's newest entry assembles the pieces above into the
presentation this library was built for.  For every $n$, the $n$-qubit
Clifford group modulo global phases is an extension of the symplectic
group by the phaseless Pauli group~\cite{selinger2015clifford}:
\[
  1 \longrightarrow (\mathbb{Z}_2 \times \mathbb{Z}_2)^n
    \longrightarrow \mathcal{C}_n
    \longrightarrow Sp(2n, \mathbb{Z}_2)
    \longrightarrow 1 ,
\]
conjugation by a Clifford operator acting on Pauli operators
symplectically --- and at $p = 2$ the extension does \emph{not} split.
The library now contains this presentation, for all $n$ at once, as an
instance of the extension recipe of \S\ref{sec:constructions}
(\texttt{Examples.Groups.Clifford.Qubit.Presentation}):

\begin{agdacode}
Clifford-pres : (n : ℕ) → WRel (PauliGen n ⊎ Gen n)
Clifford-pres n = extension-presentation (Γ-H ⊕^ n) (n QRel,_===_) conj corr
\end{agdacode}

Every argument is an artifact from earlier sections.  The normal factor
$\aid{Γ-H ⊕\char94\ n}$ is the verified Pauli presentation of
\S\ref{sec:pauli}, instantiated at $p = 2$: an $X$- and a $Z$-generator
per qubit.  The quotient relations $n\;\aid{QRel},{=}{=}{=}$ are the
symplectic relations of the staircase layer (\S\ref{sec:wip}): order and
braid-style laws for the parameterized generators $S$, $H$, and $C\!Z$,
far-commutation, and the shift-congruence \aid{cong↑}.  The action
\aid{conj} is computed rather than postulated --- a quotient generator
$g$ conjugates a Pauli generator $y$ through the symplectic action
\aid{act1} on Pauli vectors, and the result is read back into syntax:

\begin{agdacode}
conj : ∀ {n} → Gen n → PauliGen n → Word (PauliGen n)
conj g y = vecToWord (act1 g (genToVec y))
\end{agdacode}

where \aid{genToVec} reads a generator as a basis vector and
\aid{vecToWord} delogs a vector $(a_i, b_i)_i$ to the word
$\prod_i X_i^{a_i} Z_i^{b_i}$ --- the Pauli normal form of
\S\ref{sec:pauli} reused as a syntactic section.  The final ingredient
is the \emph{cocycle}, the mathematical heart of a non-split extension.
A presentation of an extension must record, for each quotient relator,
the correction: the element of the normal subgroup by which the lift of
the left-hand side differs from the lift of the right-hand side
(\S\ref{sec:constructions}).  At $p = 2$ exactly one symplectic relator
fails to lift on the nose: the phase gate obeys $S^2 = Z$, not
$S^2 = I$, on the qubit it acts on, so the relator \aid{order-S}
carries the correction $Z_0$; a shifted relator carries the shifted
correction; and every other relator lifts either exactly or up to a
global phase, which is invisible in the phaseless Pauli group.  In
code, the entire cocycle of the $n$-qubit Clifford extension is three
lines:

\begin{agdacode}
corr : ∀ {n} {u v} → (n QRel,_===_) u v → Word (PauliGen n)
corr order-S    = Z₀
corr (cong↑ r)  = shiftPauli (corr r)
corr _          = ε
\end{agdacode}

(The correction data was validated, during development, against the
exact $2{\times}2$, $4{\times}4$, and $8{\times}8$ Clifford matrices
under the same conventions.)  Figure~\ref{fig:clifford-relations}
displays the resulting relation set in full, one circuit equation per
relation family, drawn in matrix-multiplication order --- the mirror
image of the usual circuit-diagram convention, so that each diagram
transcribes its Agda word letter for letter.  A few reading notes.  In
(P4), $P$ and $Q$ range over Pauli letters on distinct qubits --- the
cross-qubit commutations of the $n$-fold product.  The conjugation
family (T1)--(T8) is shown for the standard letters; the actual family
covers every parameterized letter --- $S^k$ sends $X \mapsto X Z^k$,
the letter $H^2$ conjugates trivially at $p = 2$, and $H^3$ conjugates
as $H$ --- and every shifted letter, acting on the correspondingly
shifted qubit.  Among the twisted relators, (R9) and (R10) are
Selinger's relations (C10) and (C11) with $S^{-1} = S$ at $p = 2$;
(R14)--(R17) govern the multiplier gate $M_x$ (at $p = 2$ only $x = 1$
occurs, and $M_1 = (SH)^3$); in (R18) and (R19) the box $g$ ranges
over generators shifted past the arity of the bottom gate, with
$U \in \{H, S\}$; and (R20)--(R22) collect the parameterized letters
$S^k$, $H^k$, $C\!Z^k$ into $k$-fold products of the standard ones.  For odd primes the picture simplifies:
the extension splits, every correction is $\varepsilon$, and the recipe
degenerates to the semidirect-product presentation --- the route the
odd-prime development takes.  \aid{Clifford-pres}~$n$ is thus a finite,
explicitly given relation set for the $n$-qubit Clifford group modulo
phases, typechecked under \texttt{--safe}; what remains for the full
presentation \emph{theorem} is its semantic half --- realizing the
extension record of \S\ref{sec:constructions} over the concrete
Clifford group and supplying the two factor normal forms.  The Pauli
half is in place (\S\ref{sec:pauli}, together with its verified vector
semantics); the symplectic half is the staircase tower below.

%% The complete relation set of Clifford-pres, in circuit form.
%% Convention: diagrams are read in MATRIX-MULTIPLICATION order — the
%% rightmost gate is applied first — the inverse of the usual
%% quantum-circuit convention (and of Selinger's Figure 8).  Diagrams
%% therefore transcribe the Agda words letter for letter, left to right.
%% Each equation is wrapped in an \mbox so line breaks fall only
%% between equations.
\begin{figure}[p]
\centering
\tikzset{qz/.style={row sep={0.12cm,between origins}, column sep=0.16cm}}
\newcommand{\eqtag}[1]{{\footnotesize(#1)}}
\footnotesize

\textbf{Pauli relations}\par\smallskip

\mbox{\begin{quantikz}[qz]& \gate{X} & \gate{X} & \qw\end{quantikz}
${}={}$ \begin{quantikz}[qz]& \qw & \qw\end{quantikz}\;\eqtag{P1}}\hfill
\mbox{\begin{quantikz}[qz]& \gate{Z} & \gate{Z} & \qw\end{quantikz}
${}={}$ \begin{quantikz}[qz]& \qw & \qw\end{quantikz}\;\eqtag{P2}}\hfill
\mbox{\begin{quantikz}[qz]& \gate{X} & \gate{Z} & \qw\end{quantikz}
${}={}$ \begin{quantikz}[qz]& \gate{Z} & \gate{X} & \qw\end{quantikz}\;\eqtag{P3}}\hfill
\mbox{\begin{quantikz}[qz]& \qw & \gate{Q} & \qw \\ & \gate{P} & \qw & \qw\end{quantikz}
${}={}$ \begin{quantikz}[qz]& \gate{Q} & \qw & \qw \\ & \qw & \gate{P} & \qw\end{quantikz}\;\eqtag{P4}}

\smallskip
\textbf{Conjugation relations} (\aid{ConjRelʷ conj})\par\smallskip

\mbox{\begin{quantikz}[qz]& \gate{H} & \gate{X} & \qw\end{quantikz}
${}={}$ \begin{quantikz}[qz]& \gate{Z} & \gate{H} & \qw\end{quantikz}\;\eqtag{T1}}\hfill
\mbox{\begin{quantikz}[qz]& \gate{H} & \gate{Z} & \qw\end{quantikz}
${}={}$ \begin{quantikz}[qz]& \gate{X} & \gate{H} & \qw\end{quantikz}\;\eqtag{T2}}\hfill
\mbox{\begin{quantikz}[qz]& \gate{S} & \gate{X} & \qw\end{quantikz}
${}={}$ \begin{quantikz}[qz]& \gate{X} & \gate{Z} & \gate{S} & \qw\end{quantikz}\;\eqtag{T3}}\hfill
\mbox{\begin{quantikz}[qz]& \gate{S} & \gate{Z} & \qw\end{quantikz}
${}={}$ \begin{quantikz}[qz]& \gate{Z} & \gate{S} & \qw\end{quantikz}\;\eqtag{T4}}

\smallskip

\mbox{\begin{quantikz}[qz]& \ctrl{1} & \qw & \qw \\ & \control{} & \gate{X} & \qw\end{quantikz}
${}={}$ \begin{quantikz}[qz]& \qw & \gate{Z} & \ctrl{1} & \qw \\ & \gate{X} & \qw & \control{} & \qw\end{quantikz}\;\eqtag{T5}}\hfill
\mbox{\begin{quantikz}[qz]& \ctrl{1} & \qw & \qw \\ & \control{} & \gate{Z} & \qw\end{quantikz}
${}={}$ \begin{quantikz}[qz]& \qw & \ctrl{1} & \qw \\ & \gate{Z} & \control{} & \qw\end{quantikz}\;\eqtag{T6}}\hfill
\mbox{\begin{quantikz}[qz]& \ctrl{1} & \gate{X} & \qw \\ & \control{} & \qw & \qw\end{quantikz}
${}={}$ \begin{quantikz}[qz]& \qw & \gate{X} & \ctrl{1} & \qw \\ & \gate{Z} & \qw & \control{} & \qw\end{quantikz}\;\eqtag{T7}}\hfill
\mbox{\begin{quantikz}[qz]& \ctrl{1} & \gate{Z} & \qw \\ & \control{} & \qw & \qw\end{quantikz}
${}={}$ \begin{quantikz}[qz]& \gate{Z} & \ctrl{1} & \qw \\ & \qw & \control{} & \qw\end{quantikz}\;\eqtag{T8}}

\smallskip
\textbf{Twisted symplectic relators} (\aid{RelTwist}; (R1) is the cocycle)\par\smallskip

\mbox{\fbox{\begin{quantikz}[qz]& \gate{S} & \gate{S} & \qw\end{quantikz}
${}={}$ \begin{quantikz}[qz]& \gate{Z} & \qw\end{quantikz}}\;\eqtag{R1}}\hfill
\mbox{\begin{quantikz}[qz]& \gate{H} & \gate{H} & \gate{H} & \gate{H} & \qw\end{quantikz}
${}={}$ \begin{quantikz}[qz]& \qw & \qw\end{quantikz}\;\eqtag{R2}}\hfill
\mbox{\begin{quantikz}[qz]& \gate{S} & \gate{H} & \gate{S} & \gate{H} & \gate{S} & \gate{H} & \qw\end{quantikz}
${}={}$ \begin{quantikz}[qz]& \qw & \qw\end{quantikz}\;\eqtag{R3}}\hfill
\mbox{\begin{quantikz}[qz]& \gate{H} & \gate{H} & \gate{S} & \qw\end{quantikz}
${}={}$ \begin{quantikz}[qz]& \gate{S} & \gate{H} & \gate{H} & \qw\end{quantikz}\;\eqtag{R4}}

\smallskip

\mbox{\begin{quantikz}[qz]& \ctrl{1} & \ctrl{1} & \qw \\ & \control{} & \control{} & \qw\end{quantikz}
${}={}$ \begin{quantikz}[qz]& \qw & \qw \\ & \qw & \qw\end{quantikz}\;\eqtag{R5}}\hfill
\mbox{\begin{quantikz}[qz]& \ctrl{1} & \qw & \qw \\ & \control{} & \gate{S} & \qw\end{quantikz}
${}={}$ \begin{quantikz}[qz]& \qw & \ctrl{1} & \qw \\ & \gate{S} & \control{} & \qw\end{quantikz}\;\eqtag{R6}}\hfill
\mbox{\begin{quantikz}[qz]& \ctrl{1} & \gate{S} & \qw \\ & \control{} & \qw & \qw\end{quantikz}
${}={}$ \begin{quantikz}[qz]& \gate{S} & \ctrl{1} & \qw \\ & \qw & \control{} & \qw\end{quantikz}\;\eqtag{R7}}\hfill
\mbox{\begin{quantikz}[qz]& \ctrl{1} & \qw & \qw \\ & \control{} & \ctrl{1} & \qw \\ & \qw & \control{} & \qw\end{quantikz}
${}={}$ \begin{quantikz}[qz]& \qw & \ctrl{1} & \qw \\ & \ctrl{1} & \control{} & \qw \\ & \control{} & \qw & \qw\end{quantikz}\;\eqtag{R8}}

\smallskip

\mbox{\begin{quantikz}[qz]& \ctrl{1} & \gate{H} & \ctrl{1} & \qw \\ & \control{} & \qw & \control{} & \qw\end{quantikz}
${}={}$ \begin{quantikz}[qz]& \gate{S} & \gate{H} & \gate{S} & \ctrl{1} & \gate{H} & \gate{S} & \qw & \qw \\ & \qw & \qw & \qw & \control{} & \qw & \qw & \gate{S} & \qw\end{quantikz}\;\eqtag{R9}}\hfill
\mbox{\begin{quantikz}[qz]& \ctrl{1} & \qw & \ctrl{1} & \qw \\ & \control{} & \gate{H} & \control{} & \qw\end{quantikz}
${}={}$ \begin{quantikz}[qz]& \qw & \qw & \qw & \ctrl{1} & \qw & \qw & \gate{S} & \qw \\ & \gate{S} & \gate{H} & \gate{S} & \control{} & \gate{H} & \gate{S} & \qw & \qw\end{quantikz}\;\eqtag{R10}}

\smallskip

\mbox{\begin{quantikz}[qz]& \gate[2]{⊤⊥} & \qw & \gate[2]{⊥⊤} & \qw \\ & & \ctrl{1} & & \qw \\ & \qw & \control{} & \qw & \qw\end{quantikz}
${}={}$ \begin{quantikz}[qz]& \qw & \ctrl{1} & \qw & \qw \\ & \gate[2]{⊥⊤} & \control{} & \gate[2]{⊤⊥} & \qw \\ & & \qw & & \qw\end{quantikz}\;\eqtag{R11}}\hfill
\mbox{\begin{quantikz}[qz]& \gate[2]{⊤⊥} & \qw & \gate[2]{⊤⊥} & \qw & \gate[2]{⊤⊥} & \qw & \qw \\ & & \ctrl{1} & & \ctrl{1} & & \ctrl{1} & \qw \\ & \qw & \control{} & \qw & \control{} & \qw & \control{} & \qw\end{quantikz}
${}={}$ \begin{quantikz}[qz]& \qw & \qw \\ & \qw & \qw \\ & \qw & \qw\end{quantikz}\;\eqtag{R12}}

\smallskip

\mbox{\begin{quantikz}[qz]& \qw & \ctrl{1} & \qw & \ctrl{1} & \qw & \ctrl{1} & \qw \\ & \gate[2]{⊥⊤} & \control{} & \gate[2]{⊥⊤} & \control{} & \gate[2]{⊥⊤} & \control{} & \qw \\ & & \qw & & \qw & & \qw & \qw\end{quantikz}
${}={}$ \begin{quantikz}[qz]& \qw & \qw \\ & \qw & \qw \\ & \qw & \qw\end{quantikz}\;\eqtag{R13}}\hfill
\mbox{\begin{quantikz}[qz]& \gate{M_x} & \gate{M_y} & \qw\end{quantikz}
${}={}$ \begin{quantikz}[qz]& \gate{M_{xy}} & \qw\end{quantikz}\;\eqtag{R14}}\hfill
\mbox{\begin{quantikz}[qz]& \gate{M_x} & \gate{S} & \qw\end{quantikz}
${}={}$ \begin{quantikz}[qz]& \gate{S^{x^2}} & \gate{M_x} & \qw\end{quantikz}\;\eqtag{R15}}

\smallskip

\mbox{\begin{quantikz}[qz]& \gate{M_x} & \qw & \qw \\ & \qw & \ctrl{-1} & \qw\end{quantikz}
${}={}$ \begin{quantikz}[qz]& \gate[2]{C\!Z^{x}} & \gate{M_x} & \qw \\ & & \qw & \qw\end{quantikz}\;\eqtag{R16}}\hfill
\mbox{\begin{quantikz}[qz]& \qw & \ctrl{1} & \qw \\ & \gate{M_x} & \control{} & \qw\end{quantikz}
${}={}$ \begin{quantikz}[qz]& \gate[2]{C\!Z^{x}} & \qw & \qw \\ & & \gate{M_x} & \qw\end{quantikz}\;\eqtag{R17}}\hfill
\mbox{\begin{quantikz}[qz]& \gate{g} & \qw & \qw \\ & \qw & \gate{U} & \qw\end{quantikz}
${}={}$ \begin{quantikz}[qz]& \qw & \gate{g} & \qw \\ & \gate{U} & \qw & \qw\end{quantikz}\;\eqtag{R18}}\hfill
\mbox{\begin{quantikz}[qz]& \gate{g} & \qw & \qw \\ & \qw & \ctrl{1} & \qw \\ & \qw & \control{} & \qw\end{quantikz}
${}={}$ \begin{quantikz}[qz]& \qw & \gate{g} & \qw \\ & \ctrl{1} & \qw & \qw \\ & \control{} & \qw & \qw\end{quantikz}\;\eqtag{R19}}

\smallskip

\mbox{\begin{quantikz}[qz]& \gate{S^k} & \qw\end{quantikz}
${}={}$ \begin{quantikz}[qz]& \gate{S} & \push{\,\cdots\,} & \gate{S} & \qw\end{quantikz}\;\eqtag{R20}}\hfill
\mbox{\begin{quantikz}[qz]& \gate{H^k} & \qw\end{quantikz}
${}={}$ \begin{quantikz}[qz]& \gate{H} & \push{\,\cdots\,} & \gate{H} & \qw\end{quantikz}\;\eqtag{R21}}\hfill
\mbox{\begin{quantikz}[qz]& \gate[2]{C\!Z^{k}} & \qw \\ & & \qw\end{quantikz}
${}={}$ \begin{quantikz}[qz]& \ctrl{1} & \push{\,\cdots\,} & \ctrl{1} & \qw \\ & \control{} & \qw & \control{} & \qw\end{quantikz}\;\eqtag{R22}}

\smallskip
\textbf{Abbreviations}\par\smallskip

\mbox{\begin{quantikz}[qz]& \gate[2]{⊥⊤} & \qw \\ & & \qw\end{quantikz}
${}:={}$ \begin{quantikz}[qz]& \qw & \ctrl{1} & \qw & \gate{H} & \ctrl{1} & \gate{H} & \qw \\ & \gate{H} & \control{} & \gate{H} & \qw & \control{} & \qw & \qw\end{quantikz}}\hfill
\mbox{\begin{quantikz}[qz]& \gate[2]{⊤⊥} & \qw \\ & & \qw\end{quantikz}
${}:={}$ \begin{quantikz}[qz]& \gate{H} & \ctrl{1} & \gate{H} & \qw & \ctrl{1} & \qw & \qw \\ & \qw & \control{} & \qw & \gate{H} & \control{} & \gate{H} & \qw\end{quantikz}}\hfill
\mbox{\begin{quantikz}[qz]& \gate{M_x} & \qw\end{quantikz}
${}:={}$ \begin{quantikz}[qz]& \gate{S^x} & \gate{H} & \gate{S^{x^{-1}}} & \gate{H} & \gate{S^x} & \gate{H} & \qw\end{quantikz}}

\caption{The complete relation set of $\aid{Clifford-pres}\ n$
(\S\ref{sec:cliffordqubit}) in circuit form, at $p = 2$.
\emph{Reading convention:} diagrams compose in matrix-multiplication
order --- the \emph{rightmost} gate is applied first --- the mirror
image of the usual quantum-circuit convention (and of Figure~8
of~\cite{selinger2015clifford}), so each diagram transcribes its Agda
word letter for letter, left to right.  Wires are numbered from the
bottom; every relation may be shifted upward to any wires
(\aid{cong↑}) and is stated on the fewest wires at which it makes
sense.  The boxed (R1) is the cocycle: the lift of the quotient
relator $S^2 = \varepsilon$ picks up the Pauli correction $Z$; all
other corrections are trivial.  See \S\ref{sec:cliffordqubit} for the
conventions governing the parameterized and shifted gates.}
\label{fig:clifford-relations}
\end{figure}


\subsection{Work in progress: symplectic groups and odd-prime Clifford circuits}\label{sec:wip}

Beyond the verified catalogue, the repository contains a further
$\sim$67\,600 lines formalizing the road to multi-qupit Clifford
completeness in odd prime dimensions: modular arithmetic over
$\mathbb{Z}_p$, the symplectic groups $Sp(2n,\mathbb{Z}_p)$ and their
action on Pauli vectors, symplectic transvection lemmas, and
Selinger-style staircase normal forms for the Clifford generators
$S, H, C\!Z, \mathit{Ex}, M$ --- the formal counterpart of the qutrit
rule set of Li, Mosca, Ross, van~de~Wetering, and
Zhao~\cite{li2025qutrit} and its odd-prime generalization, with the
companion web normalizer of \S\ref{sec:tools}~\cite{qupit-tool}.  The
relations-and-action sublayer of this development typechecks --- it is
what \S\ref{sec:cliffordqubit} consumes --- but the staircase
normal-form tower is still under construction, and in keeping with our
ground rule --- nothing is claimed that does not typecheck under
\texttt{--safe} today --- we report the layer as the library's driving
next target rather than as a result.  The endgame, whose syntactic
scaffolding \S\ref{sec:cliffordqubit} has now assembled, is visible: the
Pauli presentations of this section as the normal subgroup, the
symplectic staircase as the coset tower, and the extension-presentation
theorem of \S\ref{sec:constructions} as the glue.
